% CP-VI-0031
% target_chapter: VI/35 — Proj
% source_unit_id: IMP-DL-40C58AD550-P001
% source: Downloads/theory-of-algebraic-geometry-4.tex L184-L307
% solution_provenance: SOURCE_EXPLICIT_SOLUTION

\subsection*{Problem statement}

We check the details of the construction of \(\operatorname{Proj}\) in lecture. Recall, in analogy with
\(\operatorname{Spec} A\), an affine scheme, we can define a projective scheme.

Let
\[
S=\bigoplus_{d=0}^{\infty}S_d
\]
be a graded ring. We will define \(\operatorname{Proj} S\).

\begin{enumerate}
	\item[(i)]
	We write
	\[
	S_+ = \bigoplus_{d=1}^{\infty} S_d,
	\]
	the irrelevant ideal. Define \(\operatorname{Proj} S\) to be the set of all homogeneous prime ideals
	of \(S\) not containing the irrelevant ideal.
	
	If \(I\subseteq S\) is a homogeneous ideal, let \(V(I)\) denote the set of all primes
	in \(\operatorname{Proj} S\) containing \(I\). Show these form the closed sets of a topology on
	\(\operatorname{Proj} S\).
	
	\item[(ii)]
	We define a sheaf \(\mathcal O\) on \(\operatorname{Proj} S\). For \(\mathfrak p\in \operatorname{Proj} S\), let
	\(S_{(\mathfrak p)}\) be the set of elements of the localization \(S_{\mathfrak p}\)
	which are homogeneous of degree zero, that is, a ratio of two elements of \(S\) of
	the same degree.
	
	For \(U\subseteq \operatorname{Proj} S\) open, define \(\mathcal O(U)\) to be the set of functions
	\[
	s:U\longrightarrow \coprod_{\mathfrak p\in U} S_{(\mathfrak p)}
	\]
	such that \(s(\mathfrak p)\in S_{(\mathfrak p)}\), and every point
	\(\mathfrak p\in U\) has an open neighbourhood \(V\) for which there exist
	homogeneous \(f,g\in S\) of the same degree, with
	\[
	g\notin \mathfrak q
	\]
	for all \(\mathfrak q\in V\), such that
	\[
	s(\mathfrak q)=\frac{f}{g}
	\]
	for all \(\mathfrak q\in V\).
	
	Then show:
	
	\begin{enumerate}
		\item[(a)]
		The stalk of \(\mathcal O\) at \(\mathfrak q\) is
		\[
		\mathcal O_{\mathfrak q}\cong S_{(\mathfrak q)}.
		\]
		
		\item[(b)]
		For any homogeneous \(f\in S_+\), let \(D_+(f)\) be the set of primes of
		\(\operatorname{Proj} S\) not containing \(f\). Show the sets \(D_+(f)\) cover \(\operatorname{Proj} S\),
		and for each such open set there is an isomorphism of locally ringed spaces
		\[
		\left(D_+(f),\mathcal O|_{D_+(f)}\right)
		\cong
		\operatorname{Spec} S_{(f)}.
		\]
		Here \(S_{(f)}\) denotes the subring of elements of degree \(0\) in the localization
		\(S_f\).
		
		\emph{Hint.} It is easy to define a map
		\[
		\psi:D_+(f)\longrightarrow \operatorname{Spec} S_{(f)}.
		\]
		However, constructing its inverse \(\theta\) is not so easy. Given
		\(\mathfrak q\in \operatorname{Spec} S_{(f)}\), set
		\[
		\mathfrak q_n
		:=
		\left\{
		s\in S_n :
		\frac{s^{\deg f}}{f^n}\in \mathfrak q
		\right\}.
		\]
		Show that \(\mathfrak q_n\) is closed under addition and that
		\[
		\theta(\mathfrak q)
		=
		\bigoplus_n \mathfrak q_n
		\]
		is a homogeneous prime ideal in \(D_+(f)\).
		
		\item[(c)]
		Conclude that \(\operatorname{Proj} S\) is a scheme.
		
		\item[(d)]
		Show that if \(k\) is an algebraically closed field, then the set of closed points
		of
		\[
		\operatorname{Proj} k[x_0,\dots,x_n]
		\]
		is in one-to-one correspondence with points of
		\[
		(k^{n+1}\setminus\{0\})/k^*
		\]
		with the usual action of \(k^*\) by scalar multiplication.
		
		Show that if
		\[
		I\subseteq k[x_0,\dots,x_n]
		\]
		is a homogeneous ideal, then the closed points of
		\[
		\operatorname{Proj} k[x_0,\dots,x_n]/I
		\]
		are in one-to-one correspondence with equivalence classes of points
		\[
		(a_0,\dots,a_n)\in (k^{n+1}\setminus\{0\})/k^*
		\]
		such that
		\[
		F(a_0,\dots,a_n)=0
		\]
		for all homogeneous \(F\in I\).
	\end{enumerate}
\end{enumerate}

\subsection*{Solution}

Let
\[
S=\bigoplus_{d\ge 0}S_d
\]
be a graded ring. The irrelevant ideal is
\[
S_+ = \bigoplus_{d\ge 1}S_d.
\]
We define
\[
\operatorname{Proj} S
=
\left\{
\mathfrak p\subset S :
\mathfrak p \text{ is homogeneous prime and } S_+\not\subseteq \mathfrak p
\right\}.
\]

\subsection{The topology on \texorpdfstring{$\operatorname{Proj} S$}{Proj S}}

For a homogeneous ideal \(I\subseteq S\), define
\[
V(I)
=
\left\{
\mathfrak p\in \operatorname{Proj} S : I\subseteq \mathfrak p
\right\}.
\]

We prove that these subsets form the closed sets of a topology.

First,
\[
V(0)=\operatorname{Proj} S,
\]
because every prime ideal contains \(0\).

Also,
\[
V(S_+)=\varnothing,
\]
because by definition no point of \(\operatorname{Proj} S\) contains \(S_+\).

Now let \(\{I_\alpha\}_{\alpha\in A}\) be a family of homogeneous ideals. Then
\[
\bigcap_{\alpha\in A}V(I_\alpha)
=
V\left(\sum_{\alpha\in A}I_\alpha\right).
\]
Indeed, a prime ideal contains every \(I_\alpha\) if and only if it contains their sum.

Thus arbitrary intersections of closed sets are closed.

Now let \(I,J\subseteq S\) be homogeneous ideals. Then
\[
V(I)\cup V(J)=V(IJ).
\]
Indeed, if \(\mathfrak p\supseteq I\), then \(\mathfrak p\supseteq IJ\), and similarly if
\(\mathfrak p\supseteq J\).

Conversely, suppose
\[
IJ\subseteq \mathfrak p.
\]
If \(I\not\subseteq \mathfrak p\), choose \(a\in I\) with \(a\notin \mathfrak p\). Then for every
\(b\in J\),
\[
ab\in IJ\subseteq \mathfrak p.
\]
Since \(\mathfrak p\) is prime and \(a\notin \mathfrak p\), we get
\[
b\in \mathfrak p.
\]
Hence \(J\subseteq \mathfrak p\). Therefore
\[
\mathfrak p\in V(I)\cup V(J).
\]

So finite unions of closed sets are closed.

Hence the sets
\[
V(I)
\]
for homogeneous ideals \(I\subseteq S\) define the closed sets of a topology on
\(\operatorname{Proj} S\).

\subsection{Basic open subsets}

For a homogeneous element \(f\in S\), define
\[
D_+(f)
=
\operatorname{Proj} S\setminus V((f)).
\]
Equivalently,
\[
D_+(f)
=
\left\{
\mathfrak p\in \operatorname{Proj} S : f\notin \mathfrak p
\right\}.
\]

These are the projective analogues of the affine basic open sets
\[
D(f)\subseteq \operatorname{Spec} A.
\]

They satisfy
\[
D_+(f)\cap D_+(g)=D_+(fg).
\]
Indeed,
\[
\mathfrak p\in D_+(f)\cap D_+(g)
\]
if and only if
\[
f\notin \mathfrak p
\quad\text{and}\quad
g\notin \mathfrak p.
\]
Since \(\mathfrak p\) is prime, this is equivalent to
\[
fg\notin \mathfrak p.
\]
Hence
\[
D_+(f)\cap D_+(g)=D_+(fg).
\]

The sets \(D_+(f)\), where \(f\in S_+\) is homogeneous, cover \(\operatorname{Proj} S\). Indeed, if
\[
\mathfrak p\in \operatorname{Proj} S,
\]
then
\[
S_+\not\subseteq \mathfrak p.
\]
So there exists a homogeneous element
\[
f\in S_+
\]
such that
\[
f\notin \mathfrak p.
\]
Thus
\[
\mathfrak p\in D_+(f).
\]

Therefore the \(D_+(f)\) cover \(\operatorname{Proj} S\).

\subsection{The rings \texorpdfstring{$S_{(\mathfrak p)}$}{S(p)}}

Let \(\mathfrak p\in \operatorname{Proj} S\). We define
\[
S_{(\mathfrak p)}
\]
to be the degree-zero part of the homogeneous localization of \(S\) at
\(\mathfrak p\).

Concretely,
\[
S_{(\mathfrak p)}
=
\left\{
\frac{a}{b} :
a,b\in S \text{ homogeneous},\ \deg a=\deg b,\ b\notin \mathfrak p
\right\}.
\]

This is a local ring. Its maximal ideal is
\[
\mathfrak p_{(\mathfrak p)}
=
\left\{
\frac{a}{b}\in S_{(\mathfrak p)} : a\in \mathfrak p
\right\}.
\]

Indeed, if
\[
\frac{a}{b}\notin \mathfrak p_{(\mathfrak p)},
\]
then
\[
a\notin \mathfrak p.
\]
Since also \(b\notin \mathfrak p\), the inverse
\[
\frac{b}{a}
\]
is defined in \(S_{(\mathfrak p)}\). Therefore every element outside
\(\mathfrak p_{(\mathfrak p)}\) is invertible. Hence \(S_{(\mathfrak p)}\) is local.

\subsection{The structure sheaf}

For an open subset \(U\subseteq \operatorname{Proj} S\), the ring \(\mathcal O(U)\) consists of
functions
\[
s:U\longrightarrow \coprod_{\mathfrak p\in U}S_{(\mathfrak p)}
\]
such that
\[
s(\mathfrak p)\in S_{(\mathfrak p)}
\]
for each \(\mathfrak p\in U\), and such that \(s\) is locally represented by a fraction
\[
\frac{f}{g}
\]
where \(f,g\in S\) are homogeneous of the same degree.

This is exactly the projective analogue of the structure sheaf on \(\operatorname{Spec} A\), except
that only degree-zero homogeneous fractions are allowed.

\subsection{The stalk of \texorpdfstring{$\mathcal O$}{O} at a point}

We prove that
\[
\mathcal O_{\mathfrak q}\cong S_{(\mathfrak q)}.
\]

Define
\[
\Phi:\mathcal O_{\mathfrak q}\longrightarrow S_{(\mathfrak q)}
\]
by evaluation at \(\mathfrak q\). If a germ is represented by a section \(s\in\mathcal O(U)\),
with \(\mathfrak q\in U\), define
\[
\Phi(s_{\mathfrak q})=s(\mathfrak q).
\]

This is a ring homomorphism.

We show that \(\Phi\) is surjective. Let
\[
\frac{a}{b}\in S_{(\mathfrak q)}.
\]
Then \(a,b\) are homogeneous of the same degree and
\[
b\notin \mathfrak q.
\]
Hence
\[
\mathfrak q\in D_+(b).
\]
Define a section on \(D_+(b)\) by
\[
s(\mathfrak p)=\frac{a}{b}.
\]
This is regular by definition. Its germ at \(\mathfrak q\) maps to \(a/b\). Therefore
\(\Phi\) is surjective.

Now suppose
\[
\Phi(s_{\mathfrak q})=0.
\]
Since \(s\) is locally a homogeneous fraction, after shrinking around \(\mathfrak q\) we may
write
\[
s(\mathfrak p)=\frac{a}{b}
\]
on an open neighbourhood \(V\) of \(\mathfrak q\), where \(a,b\) are homogeneous of the
same degree and \(b\notin\mathfrak p\) for all \(\mathfrak p\in V\).

The equality
\[
s(\mathfrak q)=0
\]
means
\[
\frac{a}{b}=0
\]
in \(S_{(\mathfrak q)}\). Hence there exists a homogeneous element
\[
h\notin \mathfrak q
\]
such that
\[
ha=0.
\]
Now restrict to
\[
V\cap D_+(h).
\]
On this smaller neighbourhood, \(h\notin \mathfrak p\), so
\[
\frac{a}{b}=0
\]
in \(S_{(\mathfrak p)}\) for every \(\mathfrak p\in V\cap D_+(h)\). Hence the germ
\(s_{\mathfrak q}\) is zero.

Therefore \(\Phi\) is injective and surjective. Thus
\[
\boxed{
	\mathcal O_{\mathfrak q}\cong S_{(\mathfrak q)}.
}
\]

\subsection{The affine open pieces}

Let \(f\in S_+\) be homogeneous and set
\[
d=\deg f.
\]
The localization \(S_f\) is naturally graded. Let
\[
S_{(f)}
=
(S_f)_0
\]
be its degree-zero part.

Explicitly,
\[
S_{(f)}
=
\left\{
\frac{a}{f^m} :
a\in S_{md}
\right\}.
\]

We prove that
\[
\left(D_+(f),\mathcal O|_{D_+(f)}\right)
\cong
\operatorname{Spec} S_{(f)}.
\]

\subsection{The map \texorpdfstring{$D_+(f)\to \operatorname{Spec} S_{(f)}$}{D+(f) to Spec S(f)}}

Let
\[
\mathfrak p\in D_+(f).
\]
Then \(f\notin \mathfrak p\), so we can localize \(\mathfrak p\) in \(S_f\). Define
\[
\psi(\mathfrak p)
=
\mathfrak pS_f\cap S_{(f)}.
\]

This is a prime ideal of \(S_{(f)}\), because \(\mathfrak pS_f\) is prime in \(S_f\), and
the contraction of a prime ideal is prime.

Therefore we have a map
\[
\psi:D_+(f)\longrightarrow \operatorname{Spec} S_{(f)}.
\]

\subsection{The inverse map}

Now let
\[
\mathfrak q\in \operatorname{Spec} S_{(f)}.
\]
We construct a homogeneous prime ideal of \(S\).

For \(n\ge 0\), define
\[
\mathfrak q_n
=
\left\{
s\in S_n :
\frac{s^d}{f^n}\in \mathfrak q
\right\}.
\]
This makes sense because if \(s\in S_n\), then
\[
\deg(s^d)=dn
\]
and
\[
\deg(f^n)=dn,
\]
so
\[
\frac{s^d}{f^n}
\]
has degree zero and belongs to \(S_{(f)}\).

Now set
\[
\theta(\mathfrak q)
=
\bigoplus_{n\ge 0}\mathfrak q_n.
\]

We show that \(\theta(\mathfrak q)\) is a homogeneous prime ideal.

First, by construction,
\[
\mathfrak q_n\subseteq S_n.
\]
Therefore \(\theta(\mathfrak q)\) is homogeneous.

Now we prove that each \(\mathfrak q_n\) is closed under addition. Let
\[
s,t\in \mathfrak q_n.
\]
Then
\[
\frac{s^d}{f^n}\in \mathfrak q,
\qquad
\frac{t^d}{f^n}\in \mathfrak q.
\]
We need to show that
\[
s+t\in \mathfrak q_n,
\]
that is,
\[
\frac{(s+t)^d}{f^n}\in \mathfrak q.
\]

By the binomial theorem,
\[
\frac{(s+t)^d}{f^n}
=
\sum_{i=0}^d
\binom di
\frac{s^i t^{d-i}}{f^n}.
\]
Each summand has degree zero.

For each \(i\), consider
\[
\frac{s^i t^{d-i}}{f^n}.
\]
Its \(d\)-th power is
\[
\left(\frac{s^i t^{d-i}}{f^n}\right)^d
=
\left(\frac{s^d}{f^n}\right)^i
\left(\frac{t^d}{f^n}\right)^{d-i}.
\]
Since
\[
\frac{s^d}{f^n},\frac{t^d}{f^n}\in \mathfrak q,
\]
the right-hand side lies in \(\mathfrak q\). Since \(\mathfrak q\) is prime, it is radical.
Therefore
\[
\frac{s^i t^{d-i}}{f^n}\in \mathfrak q.
\]
Hence every term in the binomial expansion lies in \(\mathfrak q\), so
\[
\frac{(s+t)^d}{f^n}\in \mathfrak q.
\]
Thus
\[
s+t\in \mathfrak q_n.
\]

Now we show that \(\theta(\mathfrak q)\) is an ideal. Let
\[
s\in \mathfrak q_n
\]
and let
\[
r\in S_m
\]
be homogeneous. Then
\[
rs\in S_{m+n}.
\]
Moreover,
\[
\frac{(rs)^d}{f^{m+n}}
=
\frac{r^d}{f^m}\cdot \frac{s^d}{f^n}.
\]
The first factor belongs to \(S_{(f)}\), and the second factor belongs to \(\mathfrak q\).
Therefore
\[
\frac{(rs)^d}{f^{m+n}}\in \mathfrak q.
\]
Hence
\[
rs\in \mathfrak q_{m+n}.
\]
So \(\theta(\mathfrak q)\) is a homogeneous ideal.

Now we prove that \(\theta(\mathfrak q)\) is prime. Let \(a\in S_m\) and \(b\in S_n\) be
homogeneous, and suppose
\[
ab\in \theta(\mathfrak q).
\]
Then
\[
ab\in \mathfrak q_{m+n}.
\]
Hence
\[
\frac{(ab)^d}{f^{m+n}}\in \mathfrak q.
\]
But
\[
\frac{(ab)^d}{f^{m+n}}
=
\frac{a^d}{f^m}\cdot \frac{b^d}{f^n}.
\]
Since \(\mathfrak q\) is prime, either
\[
\frac{a^d}{f^m}\in \mathfrak q
\]
or
\[
\frac{b^d}{f^n}\in \mathfrak q.
\]
Thus either
\[
a\in \mathfrak q_m
\]
or
\[
b\in \mathfrak q_n.
\]
Therefore either
\[
a\in \theta(\mathfrak q)
\]
or
\[
b\in \theta(\mathfrak q).
\]
So \(\theta(\mathfrak q)\) is prime.

Finally,
\[
f\notin \theta(\mathfrak q).
\]
Indeed, \(f\in S_d\), and
\[
\frac{f^d}{f^d}=1\notin \mathfrak q.
\]
So
\[
f\notin \theta(\mathfrak q).
\]
Therefore
\[
\theta(\mathfrak q)\in D_+(f).
\]

Thus we have constructed a map
\[
\theta:\operatorname{Spec} S_{(f)}\longrightarrow D_+(f).
\]

\subsection{The two maps are inverse}

We now show that
\[
\psi\circ\theta=\operatorname{id}
\qquad\text{and}\qquad
\theta\circ\psi=\operatorname{id}.
\]

Let
\[
\mathfrak q\in \operatorname{Spec} S_{(f)}.
\]
An element of \(S_{(f)}\) has the form
\[
\frac{a}{f^m},
\]
where \(a\in S_{md}\). Then
\[
\frac{a}{f^m}\in \psi(\theta(\mathfrak q))
\]
if and only if
\[
a\in \theta(\mathfrak q).
\]
By definition of \(\theta\), this is equivalent to
\[
\frac{a^d}{f^{md}}\in \mathfrak q.
\]
But
\[
\frac{a^d}{f^{md}}
=
\left(\frac{a}{f^m}\right)^d.
\]
Since \(\mathfrak q\) is prime, hence radical, this is equivalent to
\[
\frac{a}{f^m}\in \mathfrak q.
\]
Therefore
\[
\psi(\theta(\mathfrak q))=\mathfrak q.
\]

Conversely, let
\[
\mathfrak p\in D_+(f).
\]
For a homogeneous element \(s\in S_n\), we have
\[
s\in \theta(\psi(\mathfrak p))
\]
if and only if
\[
\frac{s^d}{f^n}\in \psi(\mathfrak p).
\]
This is equivalent to
\[
\frac{s^d}{f^n}\in \mathfrak pS_f.
\]
Since \(f\notin \mathfrak p\), this is equivalent to
\[
s^d\in \mathfrak p.
\]
Since \(\mathfrak p\) is prime, this is equivalent to
\[
s\in \mathfrak p.
\]
Hence
\[
\theta(\psi(\mathfrak p))=\mathfrak p.
\]

Therefore
\[
D_+(f)\cong \operatorname{Spec} S_{(f)}
\]
as topological spaces.

\subsection{Compatibility with the structure sheaves}

The homeomorphism above also identifies the structure sheaves.

Indeed, on \(D_+(f)\), sections of \(\mathcal O_{\operatorname{Proj} S}\) are locally represented by
fractions
\[
\frac{a}{b}
\]
where \(a,b\) are homogeneous of the same degree.

On \(\operatorname{Spec} S_{(f)}\), sections of the affine structure sheaf are locally represented by
fractions in localizations of \(S_{(f)}\).

These two descriptions agree under the correspondence
\[
D_+(f)\cong \operatorname{Spec} S_{(f)}.
\]

More explicitly, if \(g\in S\) is homogeneous, then
\[
D_+(f)\cap D_+(g)
\]
corresponds to the affine principal open subset
\[
D\left(\frac{g^d}{f^{\deg g}}\right)
\subseteq \operatorname{Spec} S_{(f)}.
\]
On these basic opens, both sheaves are described by the same degree-zero localized
fractions.

Therefore we obtain an isomorphism of locally ringed spaces
\[
\boxed{
	\left(D_+(f),\mathcal O|_{D_+(f)}\right)
	\cong
	\operatorname{Spec} S_{(f)}.
}
\]

\subsection{\texorpdfstring{$\operatorname{Proj} S$}{Proj S} is a scheme}

The open subsets
\[
D_+(f),
\qquad f\in S_+ \text{ homogeneous},
\]
cover \(\operatorname{Proj} S\).

For each such \(f\), we have shown that
\[
D_+(f)\cong \operatorname{Spec} S_{(f)}
\]
as locally ringed spaces.

Therefore \(\operatorname{Proj} S\) is covered by affine schemes. Hence
\[
\boxed{
	\operatorname{Proj} S \text{ is a scheme.}
}
\]
